% Figure: schematic cardiovascular branching tree with Murray's-law
% radius scaling annotations and cross-sectional area / velocity inset.
\begin{figure}[htbp]
\centering
\begin{tikzpicture}[
  axis/.style={-{Stealth}, thick, draw=engblack},
  vessel/.style={engred, line cap=round},
  label/.style={font=\scriptsize, align=left},
]

% Vessel hierarchy as horizontal segments with radius proportional to
% the diameter scale.
\foreach \i/\name/\d/\v/\area in {%
  0/aorta/0.50/0.5/4.5,
  1/large artery/0.30/0.35/15,
  2/arteriole/0.15/0.05/400,
  3/capillary/0.05/0.001/4500,
  4/venule/0.10/0.005/2500,
  5/large vein/0.30/0.15/40,
  6/vena cava/0.45/0.3/10}
{
  \pgfmathsetmacro{\xv}{2.0 * \i}
  \draw[vessel, line width=\d cm] (\xv, 3.0) -- (\xv+1.8, 3.0);
  \node[label] at (\xv+0.9, 3.6) {\name};
  \node[label, anchor=north] at (\xv+0.9, 2.85)
    {$d\sim$ \pgfmathprintnumber[fixed,precision=2]{\d}\,cm};
}

% Branching tree from aorta -> arteries -> arterioles -> capillaries
% Below the labelled hierarchy.
\begin{scope}[yshift=-1.3cm]
\node[engred, font=\bfseries] at (-0.5, 1.5) {Heart};
\draw[engred, very thick] (0, 1.5) -- (2, 1.5);

% Level 1 (large artery)
\foreach \i in {0, 1} {
  \draw[engred, thick] (2, 1.5) -- (3.5, 1.5 + 0.7*(\i - 0.5));
}
% Level 2 (arteriole)
\foreach \i in {0, 1, 2, 3} {
  \draw[engred!80, thick] (3.5, 1.85) -- (5.0, 1.85 + 0.3*(\i - 1.5));
  \draw[engred!80, thick] (3.5, 1.15) -- (5.0, 1.15 + 0.3*(\i - 1.5));
}
% Level 3 (capillary, drawn as dots beyond)
\foreach \y in {0.4, 0.7, 1.0, 1.3, 1.6, 1.9, 2.2, 2.5} {
  \draw[engred!50, thin] (5.0, \y) -- (6.5, \y);
}

\node[label, anchor=west] at (6.7, 1.5)
  {$d_{\text{p}}^3 = d_1^3 + d_2^3$\\Murray's law\\$d_{\text{daughter}}/d_{\text{parent}} = 2^{-1/3} \approx 0.794$};
\end{scope}

% Inset: cross-sectional area vs velocity
\begin{scope}[xshift=0cm, yshift=-5.5cm, scale=0.7]
\draw[axis] (0, 0) -- (10, 0) node[right, font=\small] {vessel generation};
\draw[axis] (0, 0) -- (0, 4) node[above, font=\small] {(log scale)};

% Plot total cross-sectional area (rising to capillaries, falling on return)
\draw[very thick, engblue, smooth] (0, 0.5) .. controls (1.5, 0.8) and (3.5, 2.2)
    .. (4.5, 3.5) .. controls (5.5, 2.5) and (7, 1.2) .. (9, 0.6);
\node[engblue, font=\small, anchor=south] at (4.5, 3.55) {total area};

% Velocity (falling, then rising)
\draw[very thick, engred, dashed, smooth] (0, 3.5) .. controls (1.5, 3.0) and (3.5, 1.0)
    .. (4.5, 0.5) .. controls (5.5, 1.0) and (7, 2.0) .. (9, 2.5);
\node[engred, font=\small, anchor=west] at (9.0, 2.55) {velocity};

\foreach \x/\xt in {0/aorta, 2/artery, 4/capillary, 6/venule, 8/vein} {
  \draw (\x, 0) -- (\x, -0.1) node[below, font=\tiny] {\xt};
}
\end{scope}

\end{tikzpicture}
\caption{Cardiovascular branching tree. Top: vessel hierarchy with
characteristic diameters from aorta ($\sim 1\,\text{cm}$) to capillary
($\sim 8\,\si{\micro\meter}$) and back through venules to the vena
cava. The diameter ratio at each symmetric bifurcation follows
Murray's law (the cubed-radius constraint that minimises the sum of
metabolic and pumping cost). Bottom inset: total cross-sectional area
rises by three orders of magnitude from aorta to capillary bed, and
blood velocity falls by the same factor (mass conservation in an
incompressible fluid). Slow capillary transit, $\sim 1\,\text{s}$, is
the time scale the exchange biochemistry can use.}
\label{fig:vol06ch06:cv-tree}
\end{figure}
